\section{Related Work}
\label{sec:related}

\mypara{3D understanding.}
Conventionally, deep neural architectures for understanding 3D point cloud data can be broadly classified into three categories based on their approach to modeling point clouds: projection-based, voxel-based, and point-based methods. Projection-based methods project 3D points onto various image planes and utilize 2D CNN-based backbones for feature extraction~\cite{su15mvcnn,li2016vehicle,chen2017multi,lang2019pointpillars}. Voxel-based approaches transform point clouds into regular voxel grids to facilitate 3D convolution operations~\cite{maturana2015voxnet,song2017semantic}, with their efficiency subsequently enhanced by sparse convolution~\cite{wang2017ocnn,graham20183d,choy20194d}. 
However, they often lack scalability in terms of the kernel sizes. Point-based methods, by contrast, process point clouds directly~\cite{qi2017pointnet,qi2017pointnet++,zhao2019pointweb,thomas2019kpconv, ma2022rethinking} and have recently seen a shift towards transformer-based architectures~\cite{guo2021pct,zhao2021point,wu2022point, robert2023spt, yang2023swin3d}.
While these methods are powerful, their efficiency is frequently constrained by the unstructured nature of point clouds, which poses challenges to scaling their designs.

\mypara{Serialization-based method.} 
Recent works~\cite{Wang2023OctFormer, liu2023flatformer, chen2022efficient} have introduced approaches diverging from the traditional paradigms of point cloud processing, which we categorized as serialization-based. These methods structure point clouds by sorting them according to specific patterns, transforming unstructured, irregular point clouds into manageable sequences while preserving certain spatial proximity. OctFormer~\cite{Wang2023OctFormer} inherits order during octreelization, akin to z-order, offering scalability but still constrained by the octree structure itself. FlatFormer~\cite{liu2023flatformer}, on the other hand, employs a window-based sorting strategy for grouping point pillars, akin to window partitioning. However, this design lacks scalability in the receptive field and is more suited to pillar-based 3D object detectors. These pioneering works mark the inception of serialization-based methods. Our PTv3 builds on this foundation, defining and exploring the full potential of point cloud serialization.

\mypara{3D representation learning.}
In contrast to 2D domains, where large-scale pre-training has become a standard approach for enhancing downstream tasks~\cite{caron2021emerging}, 3D representation learning is still in a phase of exploration. Most studies still rely on training models from scratch using specific target datasets~\cite{xie2020pointcontrast}. While major efforts in 3D representation learning focused on individual objects~\cite{wang2019deep,sauder2019self,sanghi2020info3d, Pang2022pointmae, Yu2022pointbert}, some recent advancements have redirected attention towards training on real-world scene-centric point clouds~\cite{xie2020pointcontrast,hou2021exploring,wu2023masked, jiang2023msp, zhu2023ponderv2}. This shift signifies a major step forward in 3D scene understanding. Notably, Point Prompt Training (PPT)~\cite{wu2023ppt} introduces a new paradigm for large-scale representation learning through multi-dataset synergistic learning, emphasizing the importance of scale. This approach greatly influences our design philosophy and initial motivation for developing PTv3, and we have incorporated this strategy in our final results.
